%!TEX root = ../main.tex

\section{Introduction}
Large language models rely on knowledge retrieval to ground and align their outputs in external evidence~\cite{ren2025investigating,wang2025survey,xia2025ground,zhang2024knowledge}, from retrieval-augmented generation (RAG) to systems and memory modules that operate over \emph{semi-structured} knowledge bases (SKB) that mix text with relational information~\cite{lewis_retrieval-augmented_2020,guu_realm_2020,karpukhin_dense_2020,izacard_leveraging_2021,mavromatis_gnn-rag_2025,chen_gril_2025,li2025retrollm}. Knowledge graphs (KGs) are a natural data representation for this setting because they organize evidence around entities and typed edges, support reuse across queries, and enforce relational constraints that a flat text index cannot express. This has motivated graph-aware retrievers and graph-grounded generation methods, including graph-based RAG and SKB retrievers that combine text and relational data~\cite{edge_local_2025,zhu_knowledge_2025,xia_knowledge-aware_2025, yao2025expandr,jeong2025database}.

Retrieving evidence from KGs is challenging because it requires coordinating two competing search modes~\cite{wu_stark_2024,lee_hybgrag_2025,zhu-etal-2025-mitigating,yan2026explore}. Many queries require \emph{breadth}: they mention multiple entities or loosely connected concepts, so the retriever must cover the graph broadly to reach the right region. Other queries require \emph{depth}: the supporting evidence only appears after following specific multi-hop relational paths. Similarity-based retrievers provide global coverage but often remain shallow and underuse relational structure, whereas traversal-based methods can be brittle because they must choose seed entities to start exploration; when seeds are incomplete or ambiguous, the search stays local and misses evidence elsewhere.

Existing work tackles these breadth-depth requirements in isolation rather than jointly~\cite{ko2025cooprags,ma2025think,wang2025corag}. \emph{Structure-aware} retrievers extend text-based retrieval with relational structure, for example, by learning node embeddings that aggregate information from nearby neighbors or by generating candidates using a local graph neighborhood before ranking them~\cite{lee_hybgrag_2025,zhu_knowledge_2025,lei_mixture_2025}. These methods capture local structure, but they typically encode a fixed neighborhood around each node, so deeper multi-hop queries require expanding the encoded context or stacking additional message-passing and retrieval stages, which increases complexity and cost~\cite{wan2025digest,hu2025remindrag,verma2024plan}. By contrast, \emph{traversal-based approaches} perform multi-hop exploration, but they depend on identifying a small set of seed entities from which exploration begins~\cite{markowitz_tree--traversals_2024,sun_think--graph_2024}. When seeds are incomplete or ambiguous, exploration stays local and misses relevant information elsewhere in the graph~\cite{liu2025ontology,ma2025think}. Many systems also rely on task- or graph-specific training to learn traversal or scoring policies, which limits transfer across domains and graph schemas~\cite{li_multi-field_2025,lei_mixture_2025,wu_avatar_2024,yu_can_2025}. As a result, existing methods struggle to combine global search with targeted relational reasoning for adaptive retrieval.

%~\cite{yao_react_2023,schick_toolformer_2023}

\textbf{Present work.}
We introduce \textbf{\retrieverShort}: \retriever, a tool-using KG retriever that gives a language model control over evidence discovery using a minimal set of tools for global lexical search and neighborhood exploration. \retrieverShort does not require selecting seed entities to start exploration or establishing a maximum hop depth in advance; instead, the model alternates between broad global search and targeted multi-hop traversal based on the query and what it has retrieved so far. We evaluate \retrieverShort on the heterogeneous graphs in the STaRK retrieval benchmark and show consistent gains across all settings. We further study compute-accuracy tradeoffs by varying the tool-call budget and the number of parallel agents, and we distill \retrieverShort's tool-use policy into an 8B model via label-free trajectory imitation, preserving most of the performance of teacher model at substantially lower inference cost~\cite{kang_distilling_2025}.

Our contributions are threefold. (i) We introduce \retriever, a training-free retrieval framework that equips language models with a minimal but expressive tool interface for adaptive retrieval from KGs. (ii) We show that \retrieverShort balances breadth-oriented retrieval with depth-oriented multi-hop traversal without task- or graph-specific training, achieving strong performance on STaRK. (iii) We distill \retrieverShort's tool-use policy without labeled supervision into a compact Qwen3-8B model~\cite{yang_qwen3_2025,kang_distilling_2025}, preserving retrieval quality while reducing inference cost.
